\section{Introduction}\label{sec:introduction}

A directed feedback vertex set of a finite digraph $G$ is a set
$F\subseteq V(G)$ for which the induced digraph $G-F$ is acyclic. Write
\[
  \DFVS(G)=\min\{\abs F:G-F\text{ is acyclic}\}.
\]
Equivalently, $F$ intersects every directed cycle. A topological order
of $G-F$ certifies this property in polynomial time. The minimum-cost
version assigns a positive rational cost to each vertex and minimizes
the sum of the deleted costs. Throughout, digraphs are finite and
loopless; two vertices may have arcs in both directions. An
$A$-approximation, for $A\ge1$, returns a feedback vertex set of
size at most $A\DFVS(G)$.

\subsection{The approximation question}

Writing $n=\abs{V(G)}$, the general $O(\log n\log\log n)$
approximation guarantee originates
in Seymour's work on fractional packing of directed circuits
\cite{Seymour1995}. Even, Naor, Schieber, and Sudan give a constructive
algorithm covering weighted directed feedback sets \cite{Even1998}.

Ordinary NP-hardness of approximation already follows from Vertex Cover,
which asks for a smallest vertex set meeting every edge of an undirected graph.
Replacing every edge of an undirected graph by its two opposite arcs
makes a vertex set a directed feedback vertex set exactly when it is a
vertex cover. Thus the earlier hardness factors below
$10\sqrt5-21$ of Dinur and Safra \cite{DinurSafra2005}, and the factors
strictly below $\sqrt2$ obtained by the completed two-to-one games line
\cite[Section~1.3.1]{KMS2018}, transfer to DFVS. These results exclude
particular constants; excluding every constant had required a stronger
complexity assumption.

Guruswami, Manokaran, and Raghavendra established
the Unique-Games barrier against every constant factor through maximum
acyclic subgraph and feedback arc set \cite{GMR2008}. The expanded
ordering-constraint framework appears in the journal paper with
H\aa stad and Charikar \cite{GHMRC2011}. Svensson gave a direct
Unique-Games-based construction with strong structure in the YES case
\cite{Svensson2013}. Guruswami and Lee subsequently gave a simpler proof
with structured completeness and short cycles in every sufficiently
large induced subgraph in the NO case \cite{GuruswamiLee2016}.
The recent discussion of Ghorbani and Mnich records the remaining gap
for general digraphs \cite{GhorbaniMnich2026}.

We prove the corresponding exclusion of every constant factor from
ordinary NP-hardness.

\begin{theorem}\label{thm:main}
For every fixed real constant $A\ge1$, it is NP-hard to approximate
minimum directed feedback vertex set within factor $A$ on unweighted
digraphs. More precisely, there is a deterministic polynomial-time
reduction from an NP-hard promise problem to instances $(G,k)$ with
$k\in\N_{>0}$. The polynomial and its constants may depend on $A$.
The reduction has the following alternatives:
\[
  \textnormal{YES: }\DFVS(G)\le k,
  \qquad
  \textnormal{NO: }\DFVS(G)>Ak.
\]
\end{theorem}

Theorem~\ref{thm:main} implies that a deterministic polynomial-time
algorithm with any constant approximation guarantee exists if and
only if $\mathrm P=\mathrm{NP}$. For the reverse implication, membership
of the size-bounded decision problem in NP permits exact optimization
and witness search when $\mathrm P=\mathrm{NP}$. The theorem establishes
that the required approximation factor grows without bound with the
input size, assuming $\mathrm P\ne\mathrm{NP}$. It does not identify
that growth rate or improve the general approximation algorithm.

A directed feedback arc set is a set of arcs whose deletion leaves an
acyclic digraph. Corollary~\ref{cor:feedback-arc} transfers the same
hardness to this problem on unweighted simple loopless digraphs, through
a reduction preserving both the optimum and approximation solutions.
The vertex-set hardness also holds when opposite arcs are forbidden,
as shown in Remark~\ref{rem:no-opposite-arcs}.

\subsection{Methodological background}

Our construction shares the useful division between protected vertices
and deletable test vertices in Svensson's dictatorship gadget
\cite[Section~3.1.1]{Svensson2013}. His soundness proof also reads a
Boolean function from a topological order
\cite[Section~3.1.3]{Svensson2013}. Here the protected vertices encode
rank profiles, and compatibility across tuple prefixes supplies a
family of monotone functions. Svensson uses the ``It Ain't Over Till
It's Over'' theorem; Guruswami and Lee instead use an invariance
principle \cite[Section~1]{GuruswamiLee2016}. Our analytic input is
Friedgut's theorem on Boolean functions of bounded total influence
\cite{Friedgut1998}. We derive its required fixed-bias form explicitly.
Junta-based label decoding already appears in Dinur and Safra's
Vertex Cover hardness proof \cite[Sections~1.1 and~3]{DinurSafra2005},
where a nearby bias is selected to obtain bounded total influence.
Here a single pivotality budget controls all the prescribed biases,
and the contradiction requires their estimates to concern the same
distribution of functions.

The rank distribution is obtained by combining finite Ramsey theory
\cite{Ramsey1930} with finite minimax \cite{vonNeumann1928}.
These are the external combinatorial tools in the rank-spreading
argument. The shared-prefix construction and the permutation
pivotality budget are proved below. Their role is to compare many
biases without changing the distribution of the functions being
analyzed. The source of hardness is the published imperfect-completeness
two-to-one theorem, so the proof does not consume a Unique Games
hardness theorem.

\subsection{The reduction and its new ingredients}

A two-to-one game has a bipartite constraint graph with a label set
$X$ at each big-side vertex and a label set $Y$ at each small-side
vertex. Each edge carries a map $X\to Y$ with exactly two preimages
per label; an edge is satisfied when the chosen endpoint labels obey
that map. The game value is the largest fraction of satisfied edges,
measured under its prescribed edge distribution.
Our starting point is the imperfect-completeness 2-to-1 Games Theorem,
obtained through the work of Dinur, Khot, Kindler, Minzer, and Safra,
Khot, Minzer, and Safra, and Barak, Kothari, and Steurer
\cite{DKKMS2025,KMS2018,BKS2019}. We use its precise form in
Theorem~\ref{thm:games}: for arbitrarily small constant $\theta>0$,
distinguishing game value at least $1-\theta$ from value at most
$\theta$ is NP-hard, with fixed alphabets and exact two-element
projection fibers. The perfect-completeness 2-to-1 Games Conjecture
is not a premise of the argument.

The central difficulty is to extract a game labeling from an arbitrary
acyclic remainder. An acyclic graph supplies a topological order, but
its order can interleave the many vertices representing the same local
test. The proof first makes comparisons within that order consistent,
then interprets them as Boolean functions. Small deletion cost will
constrain these functions in incompatible ways.

For a fixed integer factor $A$, put $D=10A$. We construct a weighted
graph with a feedback vertex set of cost at most $4$ in the YES case.
Suppose, toward the soundness contradiction, that a NO instance has
one of cost at most $D$. Fix a topological order after its deletion.
The following four steps analyze this arbitrary order. A Boolean
function on subsets is evaluated at bias $p$ by including each label
independently with probability $p$; its acceptance is the probability
that its value is one.

\begin{enumerate}[label=\textup{(\arabic*)}]
\item \emph{Rank comparisons with a common interpretation.}
For each short tuple of game vertices, we create vertices indexed by
rank-valued functions of label tuples. A global labeling evaluates
all these functions consistently. In the soundness analysis, a
putative inexpensive feedback set leaves every rank vertex present,
so its surviving topological order supplies comparison bits. A
finite Ramsey and minimax argument produces a distribution on rank
embeddings for which all triples of positions have nearly identical
image laws. Ordered comparisons then agree with high probability,
yielding a monotone Boolean function of subsets of label tuples.

\item \emph{Conditioning through a shared prefix.}
Successive random subsets split label tuples into lexicographically
ordered cells. Monotonicity selects a unique transition cell, and
prefix identifications make these transitions consistent under
refinement. A full-prefix test prevents the chosen cell from having
too few membership-one steps at any of many dyadic biases. On
averaging, the conditioned big-side Boolean function has acceptance
at least a constant multiple of the bias at every scale. The same
prefix data define functions for every game vertex before a fresh
constraint is drawn.

\item \emph{A pivotality budget independent of the alphabet size.}
A single-label test charges the event that the starred label can be
pivotal somewhere along a random permutation. The resulting quantity $B(f)$ is the sum, over labels, of the
probability that the label is pivotal at some prefix of a random
permutation of the other labels. It bounds both total influence and
acceptance probability divided by the bias, at every bias. Its
expectation is at most $D+1$ on either side of the game. Friedgut's junta theorem \cite{Friedgut1998}, transferred
to fixed dyadic biases, therefore supplies bounded coordinate lists
outside a small exceptional set. Their size is fixed before choosing
the game alphabets.

\item \emph{A common tail variable across all scales.}
A projection-respecting coupling compares bias $p$ on the big side
with bias $2p$ on the small side. Game soundness makes the two junta
lists typically use disjoint projection fibers. The corresponding
outputs are then independent under the coupling, allowing the
comparison bound to pass to product measures. At each dyadic bias, this forces a positive contribution of constant
size to the tail sum of the small-side pivotality budget. All these
estimates concern the same random variable under the same distribution.
Adding sufficiently many such contributions contradicts its expected
budget.
\end{enumerate}

The pivotality quantity is independent of the bias. This is essential:
a separate influence bound at each scale would not justify adding the
contributions to one tail sum. Likewise, the shared prefix fixes all
vertex functions before the tested game edge is drawn; otherwise the
coordinate lists would not define a game labeling.

The proof uses finite probability spaces to specify weighted test
families. Every outcome is included explicitly with its probability
as part of its cost. Thus the reduction is deterministic. All
parameters, alphabets, tuple lengths, and rank tables are constants
once the target factor is fixed, even though some are extremely
large. Finally, rounding each vertex cost upward and replacing a
vertex by an independent cluster produces an unweighted gap.

\subsection{Organization}

Section~\ref{sec:preliminaries} states the external theorems and proves
the required dyadic junta corollary.
Section~\ref{sec:construction} gives the parameter choices, graph,
tests, and completeness argument. Section~\ref{sec:ranks} proves rank
spreading and the order lemmas used to read Boolean functions from a
surviving graph. Sections~\ref{sec:prefix} and~\ref{sec:pivotality}
establish the acceptance lower bounds and pivotality budgets.
Section~\ref{sec:soundness} gives the comparison and tail-sum
contradiction. Section~\ref{sec:unweighting} removes the costs and
proves Theorem~\ref{thm:main}.
